\documentclass[12pt,a4paper]{article}

\usepackage[T2A]{fontenc}
\usepackage[utf8]{inputenc}
\usepackage[russian]{babel}
\usepackage{amsmath,amssymb,amsthm}
\usepackage{enumitem}
\usepackage{listings}
\usepackage{xcolor}
\usepackage{array}
\usepackage{longtable}
\usepackage{hyperref}

\lstdefinestyle{shell}{
    basicstyle=\ttfamily\small,
    breaklines=true,
    frame=single,
    columns=fullflexible,
    keepspaces=true,
    backgroundcolor=\color{gray!5},
    literate=
      {А}{{\CYRA}}1
      {Б}{{\CYRB}}1
      {В}{{\CYRV}}1
      {Г}{{\CYRG}}1
      {Д}{{\CYRD}}1
      {Е}{{\CYRE}}1
      {Ё}{{\CYRYO}}1
      {Ж}{{\CYRZH}}1
      {З}{{\CYRZ}}1
      {И}{{\CYRI}}1
      {Й}{{\CYRISHRT}}1
      {К}{{\CYRK}}1
      {Л}{{\CYRL}}1
      {М}{{\CYRM}}1
      {Н}{{\CYRN}}1
      {О}{{\CYRO}}1
      {П}{{\CYRP}}1
      {Р}{{\CYRR}}1
      {С}{{\CYRS}}1
      {Т}{{\CYRT}}1
      {У}{{\CYRU}}1
      {Ф}{{\CYRF}}1
      {Х}{{\CYRH}}1
      {Ц}{{\CYRC}}1
      {Ч}{{\CYRCH}}1
      {Ш}{{\CYRSH}}1
      {Щ}{{\CYRSHCH}}1
      {Ъ}{{\CYRHRDSN}}1
      {Ы}{{\CYRERY}}1
      {Ь}{{\CYRSFTSN}}1
      {Э}{{\CYREREV}}1
      {Ю}{{\CYRYU}}1
      {Я}{{\CYRYA}}1
      {а}{{\cyra}}1
      {б}{{\cyrb}}1
      {в}{{\cyrv}}1
      {г}{{\cyrg}}1
      {д}{{\cyrd}}1
      {е}{{\cyre}}1
      {ё}{{\cyryo}}1
      {ж}{{\cyrzh}}1
      {з}{{\cyrz}}1
      {и}{{\cyri}}1
      {й}{{\cyrishrt}}1
      {к}{{\cyrk}}1
      {л}{{\cyrl}}1
      {м}{{\cyrm}}1
      {н}{{\cyrn}}1
      {о}{{\cyro}}1
      {п}{{\cyrp}}1
      {р}{{\cyrr}}1
      {с}{{\cyrs}}1
      {т}{{\cyrt}}1
      {у}{{\cyru}}1
      {ф}{{\cyrf}}1
      {х}{{\cyrh}}1
      {ц}{{\cyrc}}1
      {ч}{{\cyrch}}1
      {ш}{{\cyrsh}}1
      {щ}{{\cyrshch}}1
      {ъ}{{\cyrhrdsn}}1
      {ы}{{\cyrery}}1
      {ь}{{\cyrsftsn}}1
      {э}{{\cyrerev}}1
      {ю}{{\cyryu}}1
      {я}{{\cyrya}}1
}

\lstset{style=shell}

\title{Билет №51 \\ \small Командные оболочки, языки сценариев командных оболочек. Командные оболочки различных семейств операционных систем (Windows, Unix-подобных). (СпБГУ)}
\author{Сериков Владислав Александрович}
\date{16 мая 2026}

\begin{document}

\maketitle
\tableofcontents
\newpage

\section{Понятие командной оболочки}

\textbf{Командная оболочка} --- это программа-посредник между пользователем и операционной системой, предназначенная для ввода, интерпретации и выполнения команд.

Оболочка обычно выполняет следующие функции:
\begin{enumerate}
    \item принимает команды пользователя в интерактивном режиме;
    \item читает команды из файлов сценариев;
    \item выполняет лексический и синтаксический анализ командной строки;
    \item запускает внешние программы;
    \item предоставляет встроенные команды;
    \item управляет переменными, окружением, потоками ввода-вывода;
    \item поддерживает перенаправление, конвейеры, подстановки, шаблоны имён файлов;
    \item обеспечивает средства автоматизации административных и пользовательских задач.
\end{enumerate}

Командная оболочка не является ядром операционной системы. Она использует системные вызовы или системные API для запуска процессов, работы с файлами, потоками, переменными окружения и другими ресурсами.

\subsection{Интерактивная и неинтерактивная работа}

Оболочка может работать в двух основных режимах.

\begin{description}
    \item[Интерактивный режим] Пользователь вводит команды с клавиатуры, оболочка немедленно их интерпретирует и выводит результат.
    \item[Пакетный режим] Команды считываются из файла сценария и выполняются последовательно или согласно управляющим конструкциям.
\end{description}

Пример интерактивной работы в Unix-подобной системе:

\begin{lstlisting}
$ pwd
/home/student
$ ls
main.c  notes.txt  script.sh
\end{lstlisting}

Пример сценария:

\begin{lstlisting}
#!/bin/sh
echo "Текущий каталог:"
pwd
\end{lstlisting}

\section{Командная строка и модель выполнения команды}

\subsection{Команда}

\textbf{Команда} --- это инструкция, передаваемая оболочке. Она может быть:
\begin{enumerate}
    \item встроенной командой оболочки;
    \item внешней программой;
    \item функцией оболочки;
    \item псевдонимом;
    \item управляющей конструкцией языка сценариев.
\end{enumerate}

Например, в Unix-подобной системе:

\begin{lstlisting}
echo Hello
ls -l /tmp
cd /home/student
\end{lstlisting}

Здесь \texttt{echo} может быть встроенной командой, \texttt{ls} обычно является внешней программой, а \texttt{cd} почти всегда является встроенной командой, так как изменение текущего каталога должно происходить внутри процесса оболочки.

\subsection{Общий алгоритм обработки командной строки}

Типичная командная оболочка выполняет команду в несколько этапов.

\textbf{Алгоритм обработки команды:}
\begin{enumerate}
    \item Считать строку команды.
    \item Разбить строку на лексемы с учётом кавычек, экранирования и специальных символов.
    \item Выполнить подстановки:
    \begin{enumerate}
        \item подстановку переменных;
        \item командную подстановку;
        \item арифметическую подстановку;
        \item подстановку имён файлов по шаблону;
        \item разбиение на слова, если оно предусмотрено языком оболочки.
    \end{enumerate}
    \item Распознать перенаправления ввода-вывода.
    \item Определить тип команды: встроенная команда, функция, внешняя программа, управляющая конструкция.
    \item Подготовить окружение для выполнения.
    \item Выполнить команду.
    \item Получить код завершения.
    \item Вернуть управление пользователю или перейти к следующей команде сценария.
\end{enumerate}

Пример:

\begin{lstlisting}
echo "Current directory: $(pwd)"
\end{lstlisting}

Иллюстрация шагов:
\begin{enumerate}
    \item оболочка читает строку \texttt{echo "Current directory: \$(pwd)"};
    \item распознаёт командную подстановку \texttt{\$(pwd)};
    \item выполняет команду \texttt{pwd};
    \item подставляет её результат;
    \item запускает \texttt{echo} с готовым аргументом.
\end{enumerate}

\section{Языки сценариев командных оболочек}

\textbf{Язык сценариев командной оболочки} --- это язык программирования, встроенный в оболочку или тесно связанный с ней, предназначенный для автоматизации выполнения команд операционной системы.

Обычно такой язык поддерживает:
\begin{enumerate}
    \item переменные;
    \item условные операторы;
    \item циклы;
    \item функции или процедуры;
    \item обработку аргументов командной строки;
    \item работу с кодами завершения;
    \item перенаправление потоков;
    \item конвейеры;
    \item вызов внешних программ.
\end{enumerate}

\subsection{Особенности языков сценариев оболочек}

Языки оболочек отличаются от универсальных языков программирования следующими свойствами:
\begin{enumerate}
    \item основной объект выполнения --- команда операционной системы;
    \item широко используется текстовый ввод-вывод;
    \item программы часто строятся как композиции внешних утилит;
    \item синтаксис ориентирован на краткость интерактивного ввода;
    \item ошибки часто связаны с пробелами, кавычками и подстановками;
    \item большое значение имеют коды завершения команд.
\end{enumerate}

\section{Коды завершения}

Каждая команда обычно возвращает \textbf{код завершения}. В Unix-подобных системах код \texttt{0} означает успешное завершение, а ненулевое значение --- ошибку или особое состояние.

Пример:

\begin{lstlisting}
#!/bin/sh

mkdir testdir

if [ $? -eq 0 ]; then
    echo "Каталог создан"
else
    echo "Ошибка создания каталога"
fi
\end{lstlisting}

Более идиоматический вариант:

\begin{lstlisting}
#!/bin/sh

if mkdir testdir; then
    echo "Каталог создан"
else
    echo "Ошибка создания каталога"
fi
\end{lstlisting}

В PowerShell также существует понятие успешности выполнения, но оно устроено иначе: есть автоматическая переменная \texttt{\$?}, переменная \texttt{\$LASTEXITCODE} для внешних программ и система исключений.

Пример PowerShell:

\begin{lstlisting}
New-Item -ItemType Directory -Path testdir

if ($?) {
    Write-Output "Каталог создан"
} else {
    Write-Output "Ошибка создания каталога"
}
\end{lstlisting}

\section{Потоки ввода-вывода}

\subsection{Стандартные потоки}

В Unix-подобных системах процесс обычно имеет три стандартных потока:

\begin{center}
\begin{tabular}{|c|c|c|}
\hline
Номер & Имя & Назначение \\
\hline
0 & stdin & стандартный ввод \\
1 & stdout & стандартный вывод \\
2 & stderr & стандартный поток ошибок \\
\hline
\end{tabular}
\end{center}

\subsection{Перенаправление}

Перенаправление позволяет изменить источник ввода или назначение вывода.

Примеры для Unix-подобных оболочек:

\begin{lstlisting}
ls > files.txt          # записать stdout в файл
ls >> files.txt         # дописать stdout в файл
cat < input.txt         # читать stdin из файла
ls 2> errors.txt        # записать stderr в файл
ls > out.txt 2> err.txt # разделить stdout и stderr
ls > all.txt 2>&1       # объединить stdout и stderr
\end{lstlisting}

Примеры для Windows CMD:

\begin{lstlisting}
dir > files.txt
dir >> files.txt
dir 2> errors.txt
dir > all.txt 2>&1
\end{lstlisting}

В PowerShell есть несколько потоков: success output, error, warning, verbose, debug, information. Поэтому перенаправление устроено богаче:

\begin{lstlisting}
Get-ChildItem > files.txt
Get-ChildItem 2> errors.txt
Get-ChildItem *> all.txt
\end{lstlisting}

\section{Конвейеры}

\textbf{Конвейер} --- это механизм соединения вывода одной команды со вводом другой.

В Unix-подобных оболочках конвейер обычно передаёт байтовый или текстовый поток:

\begin{lstlisting}
cat access.log | grep "ERROR" | wc -l
\end{lstlisting}

Эта команда:
\begin{enumerate}
    \item читает файл \texttt{access.log};
    \item выбирает строки, содержащие \texttt{ERROR};
    \item считает количество таких строк.
\end{enumerate}

Более экономичный вариант без лишнего \texttt{cat}:

\begin{lstlisting}
grep "ERROR" access.log | wc -l
\end{lstlisting}

В PowerShell конвейер передаёт не просто текст, а объекты .NET:

\begin{lstlisting}
Get-Process | Where-Object { $_.CPU -gt 10 } | Sort-Object CPU
\end{lstlisting}

Здесь:
\begin{enumerate}
    \item \texttt{Get-Process} возвращает объекты процессов;
    \item \texttt{Where-Object} фильтрует объекты;
    \item \texttt{Sort-Object} сортирует их по свойству \texttt{CPU}.
\end{enumerate}

Это фундаментальное отличие PowerShell от традиционных текстовых оболочек Unix.

\section{Переменные и окружение}

\subsection{Переменные оболочки}

Переменная оболочки хранит значение, доступное внутри текущей оболочки.

Пример POSIX shell / Bash:

\begin{lstlisting}
name="Alice"
echo "$name"
\end{lstlisting}

Важно: вокруг знака \texttt{=} не должно быть пробелов.

Неверно:

\begin{lstlisting}
name = "Alice"
\end{lstlisting}

\subsection{Переменные окружения}

\textbf{Переменные окружения} --- это пары ``имя--значение'', передаваемые дочерним процессам.

Пример:

\begin{lstlisting}
export PATH="/usr/local/bin:$PATH"
\end{lstlisting}

Переменная \texttt{PATH} содержит список каталогов, в которых оболочка ищет исполняемые файлы.

Минимальный пример:

\begin{lstlisting}
#!/bin/sh

MYVAR="local"
export EXPORTED="visible"

sh -c 'echo "MYVAR=$MYVAR"; echo "EXPORTED=$EXPORTED"'
\end{lstlisting}

Ожидаемый результат:

\begin{lstlisting}
MYVAR=
EXPORTED=visible
\end{lstlisting}

Переменная \texttt{MYVAR} не была экспортирована, поэтому дочерняя оболочка её не получила.

\subsection{Переменные в Windows CMD}

В CMD переменные окружения записываются через \texttt{\%NAME\%}:

\begin{lstlisting}
set NAME=Alice
echo %NAME%
\end{lstlisting}

Для изменения \texttt{PATH}:

\begin{lstlisting}
set PATH=C:\Tools;%PATH%
\end{lstlisting}

\subsection{Переменные в PowerShell}

В PowerShell переменные имеют префикс \texttt{\$}:

\begin{lstlisting}
$name = "Alice"
Write-Output $name
\end{lstlisting}

Переменные окружения доступны через псевдодиск \texttt{Env:}:

\begin{lstlisting}
$env:PATH
$env:MYVAR = "value"
\end{lstlisting}

\section{Кавычки и экранирование}

Кавычки определяют, как оболочка интерпретирует специальные символы.

\subsection{POSIX shell и Bash}

Одинарные кавычки запрещают почти все подстановки:

\begin{lstlisting}
name="Alice"
echo '$name'
\end{lstlisting}

Результат:

\begin{lstlisting}
$name
\end{lstlisting}

Двойные кавычки допускают подстановку переменных:

\begin{lstlisting}
name="Alice"
echo "$name"
\end{lstlisting}

Результат:

\begin{lstlisting}
Alice
\end{lstlisting}

Обратная косая черта экранирует следующий символ:

\begin{lstlisting}
echo "Цена: \$100"
\end{lstlisting}

Результат:

\begin{lstlisting}
Цена: $100
\end{lstlisting}

\subsection{CMD}

В CMD переменные чаще всего раскрываются через \texttt{\%VAR\%}:

\begin{lstlisting}
set NAME=Alice
echo %NAME%
\end{lstlisting}

Экранирование специальных символов часто выполняется символом \texttt{\^}:

\begin{lstlisting}
echo 1 ^> 2
\end{lstlisting}

\subsection{PowerShell}

В PowerShell одинарные кавычки задают литеральную строку:

\begin{lstlisting}
$name = "Alice"
'$name'
\end{lstlisting}

Результат:

\begin{lstlisting}
$name
\end{lstlisting}

Двойные кавычки выполняют интерполяцию:

\begin{lstlisting}
$name = "Alice"
"$name"
\end{lstlisting}

Результат:

\begin{lstlisting}
Alice
\end{lstlisting}

Для экранирования используется обратный апостроф:

\begin{lstlisting}
"Цена: `$100"
\end{lstlisting}

\section{Шаблоны имён файлов}

В Unix-подобных оболочках часто используется \textbf{globbing} --- подстановка имён файлов по шаблонам.

\begin{center}
\begin{tabular}{|c|p{9cm}|}
\hline
Шаблон & Значение \\
\hline
\texttt{*} & любая строка символов \\
\texttt{?} & ровно один символ \\
\texttt{[abc]} & один символ из множества \texttt{a}, \texttt{b}, \texttt{c} \\
\texttt{[0-9]} & одна цифра \\
\hline
\end{tabular}
\end{center}

Примеры:

\begin{lstlisting}
ls *.txt
rm file?.log
cp image[0-9].png backup/
\end{lstlisting}

Важно понимать, что шаблоны раскрывает сама оболочка до запуска программы. Например, команда:

\begin{lstlisting}
echo *.txt
\end{lstlisting}

может быть преобразована оболочкой в:

\begin{lstlisting}
echo a.txt b.txt c.txt
\end{lstlisting}

\section{Командные оболочки Unix-подобных систем}

\subsection{Исторические и современные оболочки}

В Unix-подобных системах существует множество командных оболочек.

\begin{longtable}{|p{3cm}|p{4cm}|p{7cm}|}
\hline
Оболочка & Семейство & Особенности \\
\hline
\texttt{sh} & Bourne shell & классическая Unix-оболочка, основа POSIX shell \\
\hline
\texttt{bash} & Bourne-like & GNU Bourne Again Shell, распространена в Linux \\
\hline
\texttt{dash} & POSIX shell & компактная и быстрая оболочка, часто используется как \texttt{/bin/sh} \\
\hline
\texttt{ksh} & Korn shell & расширения над Bourne shell, сильные средства сценариев \\
\hline
\texttt{zsh} & Bourne-like & развитая интерактивная оболочка, автодополнение, темы, плагины \\
\hline
\texttt{csh}, \texttt{tcsh} & C shell & синтаксис, частично напоминающий C; исторически популярен в BSD-системах \\
\hline
\texttt{fish} & Friendly Interactive Shell & удобство интерактивной работы, автодополнение, подсветка, не является POSIX-совместимой \\
\hline
\end{longtable}

\subsection{POSIX shell}

\textbf{POSIX shell} --- стандартный командный язык, определённый спецификацией POSIX для утилиты \texttt{sh}. Он задаёт переносимый минимум возможностей оболочки.

Пример POSIX-совместимого сценария:

\begin{lstlisting}
#!/bin/sh

for file in *.txt
do
    echo "File: $file"
done
\end{lstlisting}

Если сценарий должен быть переносимым между Unix-подобными системами, следует избегать специфических расширений Bash, Zsh или Fish.

\subsection{Bash}

\textbf{Bash} --- одна из наиболее распространённых Unix-подобных оболочек. Она совместима с Bourne shell в значительной степени и содержит многочисленные расширения.

К возможностям Bash относятся:
\begin{enumerate}
    \item история команд;
    \item автодополнение;
    \item массивы;
    \item арифметические выражения;
    \item функции;
    \item расширенные шаблоны;
    \item командная подстановка;
    \item процессная подстановка;
    \item управление заданиями.
\end{enumerate}

Пример Bash-скрипта с массивом:

\begin{lstlisting}
#!/usr/bin/env bash

names=("Alice" "Bob" "Charlie")

for name in "${names[@]}"; do
    echo "Hello, $name"
done
\end{lstlisting}

Такой сценарий не является строго POSIX-совместимым, поскольку массивы вида \texttt{name=(...)} не входят в POSIX shell.

\subsection{Запуск сценариев в Unix-подобных системах}

Файл сценария обычно начинается со строки \textbf{shebang}:

\begin{lstlisting}
#!/bin/sh
\end{lstlisting}

или:

\begin{lstlisting}
#!/usr/bin/env bash
\end{lstlisting}

Эта строка указывает системе, какой интерпретатор использовать для выполнения файла.

Алгоритм запуска сценария:
\begin{enumerate}
    \item Пользователь вводит имя файла сценария.
    \item Ядро проверяет право на исполнение.
    \item Если файл начинается с \texttt{\#!}, ядро извлекает путь к интерпретатору.
    \item Запускается интерпретатор.
    \item Интерпретатор читает файл сценария и выполняет команды.
\end{enumerate}

Пример:

\begin{lstlisting}
chmod +x script.sh
./script.sh
\end{lstlisting}

\section{Основы сценариев POSIX shell и Bash}

\subsection{Условный оператор}

Общий вид:

\begin{lstlisting}
if command; then
    commands
elif other_command; then
    other_commands
else
    fallback_commands
fi
\end{lstlisting}

Пример:

\begin{lstlisting}
#!/bin/sh

if [ -f "input.txt" ]; then
    echo "Файл существует"
else
    echo "Файл не найден"
fi
\end{lstlisting}

Здесь \texttt{[} --- это команда проверки условия, также известная как \texttt{test}. Между \texttt{[}, аргументами и \texttt{]} обязательны пробелы.

\subsection{Цикл for}

\begin{lstlisting}
#!/bin/sh

for file in *.txt
do
    echo "Обрабатывается файл $file"
done
\end{lstlisting}

\subsection{Цикл while}

\begin{lstlisting}
#!/bin/sh

i=1
while [ "$i" -le 5 ]
do
    echo "$i"
    i=$((i + 1))
done
\end{lstlisting}

\subsection{Чтение строк файла}

\begin{lstlisting}
#!/bin/sh

while IFS= read -r line
do
    echo "Строка: $line"
done < input.txt
\end{lstlisting}

Здесь:
\begin{enumerate}
    \item \texttt{IFS=} предотвращает удаление начальных и конечных разделителей;
    \item \texttt{read -r} запрещает специальную обработку обратной косой черты;
    \item \texttt{< input.txt} перенаправляет файл на стандартный ввод цикла.
\end{enumerate}

\subsection{Функции}

\begin{lstlisting}
#!/bin/sh

greet() {
    echo "Hello, $1"
}

greet "Alice"
greet "Bob"
\end{lstlisting}

Аргументы функции доступны как \texttt{\$1}, \texttt{\$2}, и так далее.

\subsection{Аргументы командной строки}

\begin{center}
\begin{tabular}{|c|p{9cm}|}
\hline
Переменная & Значение \\
\hline
\texttt{\$0} & имя сценария \\
\texttt{\$1}, \texttt{\$2}, \ldots & позиционные аргументы \\
\texttt{\$\#} & количество аргументов \\
\texttt{\$@} & список всех аргументов \\
\texttt{\$?} & код завершения последней команды \\
\texttt{\$\$} & идентификатор текущего процесса оболочки \\
\hline
\end{tabular}
\end{center}

Пример:

\begin{lstlisting}
#!/bin/sh

echo "Имя сценария: $0"
echo "Первый аргумент: $1"
echo "Количество аргументов: $#"
\end{lstlisting}

\section{Командные оболочки Windows}

\subsection{Общие сведения}

В Windows исторически применялись и применяются несколько средств командной работы:
\begin{enumerate}
    \item \texttt{COMMAND.COM} в MS-DOS и ранних Windows;
    \item \texttt{cmd.exe}, или Windows Command Prompt;
    \item Windows Script Host;
    \item Windows PowerShell;
    \item PowerShell 7 и новее;
    \item Windows Terminal как современный терминальный эмулятор, способный запускать разные оболочки.
\end{enumerate}

Важно различать терминал и оболочку. Терминал предоставляет окно и средства ввода-вывода, а оболочка интерпретирует команды. Например, Windows Terminal может запускать PowerShell, CMD, WSL Bash и другие оболочки.

\subsection{CMD}

\textbf{CMD} --- традиционная командная оболочка Windows NT-семейства. Она поддерживает интерактивные команды и пакетные файлы с расширениями \texttt{.bat} и \texttt{.cmd}.

Пример пакетного файла:

\begin{lstlisting}
@echo off
echo Hello, %USERNAME%
dir
\end{lstlisting}

Здесь:
\begin{enumerate}
    \item \texttt{@echo off} отключает вывод самих команд;
    \item \texttt{\%USERNAME\%} раскрывается в имя текущего пользователя;
    \item \texttt{dir} выводит содержимое каталога.
\end{enumerate}

\subsection{Условные операторы CMD}

\begin{lstlisting}
@echo off

if exist input.txt (
    echo Файл найден
) else (
    echo Файл не найден
)
\end{lstlisting}

Проверка кода завершения:

\begin{lstlisting}
@echo off

mkdir testdir

if errorlevel 1 (
    echo Ошибка
) else (
    echo Каталог создан
)
\end{lstlisting}

Особенность \texttt{if errorlevel N}: условие истинно, если код завершения больше или равен \texttt{N}.

\subsection{Циклы CMD}

Пример цикла по файлам:

\begin{lstlisting}
@echo off

for %%f in (*.txt) do (
    echo Файл: %%f
)
\end{lstlisting}

В интерактивной командной строке используется один знак процента:

\begin{lstlisting}
for %f in (*.txt) do echo %f
\end{lstlisting}

В пакетном файле используется двойной знак процента:

\begin{lstlisting}
for %%f in (*.txt) do echo %%f
\end{lstlisting}

\subsection{PowerShell}

\textbf{PowerShell} --- командная оболочка и язык сценариев, основанный на объектной модели .NET. В отличие от традиционных оболочек, PowerShell передаёт по конвейеру объекты, а не только текст.

Команды PowerShell называются \textbf{cmdlet}. Их имена обычно имеют форму:

\[
\texttt{Verb-Noun}
\]

Примеры:

\begin{lstlisting}
Get-Process
Get-Service
Set-Location
New-Item
Remove-Item
\end{lstlisting}

\subsection{Сценарии PowerShell}

Файлы сценариев PowerShell имеют расширение \texttt{.ps1}.

Пример:

\begin{lstlisting}
$name = "Alice"
Write-Output "Hello, $name"
\end{lstlisting}

Условный оператор:

\begin{lstlisting}
if (Test-Path "input.txt") {
    Write-Output "Файл найден"
} else {
    Write-Output "Файл не найден"
}
\end{lstlisting}

Цикл:

\begin{lstlisting}
foreach ($file in Get-ChildItem *.txt) {
    Write-Output "Файл: $($file.Name)"
}
\end{lstlisting}

Функция:

\begin{lstlisting}
function Say-Hello {
    param(
        [string]$Name
    )

    Write-Output "Hello, $Name"
}

Say-Hello -Name "Alice"
\end{lstlisting}

\subsection{Объектный конвейер PowerShell}

Рассмотрим команду:

\begin{lstlisting}
Get-Process | Where-Object { $_.WorkingSet -gt 100MB } |
    Sort-Object WorkingSet -Descending |
    Select-Object -First 5 Name, WorkingSet
\end{lstlisting}

Шаги выполнения:
\begin{enumerate}
    \item \texttt{Get-Process} получает список процессов как объекты.
    \item \texttt{Where-Object} оставляет только процессы, использующие более 100 МБ памяти.
    \item \texttt{Sort-Object} сортирует объекты по рабочему набору памяти.
    \item \texttt{Select-Object} выбирает первые пять объектов и нужные свойства.
\end{enumerate}

В Unix-подобной оболочке аналогичная задача часто решается обработкой текста:

\begin{lstlisting}
ps aux | sort -nrk 6 | head -5
\end{lstlisting}

Однако текстовый подход менее структурирован: программа должна знать, в каком столбце находится нужное значение.

\section{Сравнение Unix shell, CMD и PowerShell}

\begin{longtable}{|p{4cm}|p{3.5cm}|p{3.5cm}|p{3.5cm}|}
\hline
Критерий & Unix shell / Bash & CMD & PowerShell \\
\hline
Основная модель & текстовые команды и потоки & текстовые команды Windows & объектный конвейер .NET \\
\hline
Тип сценариев & \texttt{.sh} & \texttt{.bat}, \texttt{.cmd} & \texttt{.ps1} \\
\hline
Переменные & \texttt{\$VAR} & \texttt{\%VAR\%} & \texttt{\$var} \\
\hline
Переменные окружения & \texttt{\$PATH}, \texttt{export} & \texttt{\%PATH\%}, \texttt{set} & \texttt{\$env:PATH} \\
\hline
Конвейер & байты/текст & текст & объекты \\
\hline
Условие успеха & код завершения 0 & \texttt{ERRORLEVEL} & \texttt{\$?}, исключения, \texttt{\$LASTEXITCODE} \\
\hline
Типичные задачи & администрирование Unix/Linux, автоматизация, сборка & автоматизация Windows-задач старого типа & современное администрирование Windows и кроссплатформенная автоматизация \\
\hline
\end{longtable}

\section{Встроенные и внешние команды}

\textbf{Встроенная команда} выполняется самой оболочкой. \textbf{Внешняя команда} является отдельной программой.

Примеры встроенных команд POSIX shell:
\begin{enumerate}
    \item \texttt{cd};
    \item \texttt{echo};
    \item \texttt{export};
    \item \texttt{read};
    \item \texttt{set};
    \item \texttt{alias} в оболочках, где он поддерживается.
\end{enumerate}

Команда \texttt{cd} должна быть встроенной, потому что текущий каталог является свойством процесса. Если бы \texttt{cd} выполнялась внешней программой, она изменила бы каталог только своего дочернего процесса, после чего этот процесс завершился бы, а каталог родительской оболочки остался бы прежним.

Иллюстрация:

\begin{lstlisting}
# Внешняя программа не может изменить текущий каталог родительской оболочки
some_external_cd /tmp
pwd
\end{lstlisting}

Поэтому \texttt{cd} реализуется внутри оболочки.

\section{Поиск исполняемых файлов}

Когда пользователь вводит команду, оболочка должна определить, что именно выполнять.

\subsection{Алгоритм поиска команды в Unix-подобной оболочке}

Упрощённый алгоритм:
\begin{enumerate}
    \item Если команда содержит символ \texttt{/}, рассматривать её как путь.
    \item Иначе проверить, не является ли команда ключевым словом языка оболочки.
    \item Проверить, не является ли команда функцией.
    \item Проверить, не является ли команда встроенной командой.
    \item Проверить псевдонимы, если они применимы в данном контексте.
    \item Последовательно просмотреть каталоги из переменной \texttt{PATH}.
    \item Если исполняемый файл найден, запустить его.
    \item Если не найден, вывести ошибку.
\end{enumerate}

Пример:

\begin{lstlisting}
echo $PATH
which ls
command -v ls
\end{lstlisting}

\subsection{Поиск команд в CMD}

CMD ищет исполняемые файлы с учётом переменной \texttt{PATH} и расширений из \texttt{PATHEXT}.

Например, если введено:

\begin{lstlisting}
notepad
\end{lstlisting}

CMD может найти файл:

\begin{lstlisting}
C:\Windows\System32\notepad.exe
\end{lstlisting}

\section{Управление заданиями в Unix-подобных оболочках}

\textbf{Задание} --- это команда или конвейер, запущенный оболочкой и отслеживаемый ею.

Запуск в фоновом режиме:

\begin{lstlisting}
sleep 100 &
\end{lstlisting}

Просмотр заданий:

\begin{lstlisting}
jobs
\end{lstlisting}

Возврат задания на передний план:

\begin{lstlisting}
fg %1
\end{lstlisting}

Продолжение в фоне:

\begin{lstlisting}
bg %1
\end{lstlisting}

Завершение процесса:

\begin{lstlisting}
kill %1
\end{lstlisting}

Управление заданиями тесно связано с процессами, группами процессов, терминалом и сигналами.

\section{Сигналы и обработка завершения}

В Unix-подобных системах процессы могут получать сигналы. Например:
\begin{enumerate}
    \item \texttt{SIGINT} --- обычно посылается при нажатии \texttt{Ctrl+C};
    \item \texttt{SIGTERM} --- вежливая просьба завершиться;
    \item \texttt{SIGKILL} --- принудительное завершение;
    \item \texttt{SIGHUP} --- разрыв соединения с терминалом или сигнал перечитать конфигурацию.
\end{enumerate}

В shell-сценариях можно использовать \texttt{trap}.

Пример:

\begin{lstlisting}
#!/bin/sh

cleanup() {
    echo "Очистка временных файлов"
    rm -f /tmp/my_temp_file
}

trap cleanup EXIT INT TERM

echo "temporary data" > /tmp/my_temp_file
sleep 10
\end{lstlisting}

Здесь функция \texttt{cleanup} будет вызвана при завершении сценария или получении некоторых сигналов.

\section{Безопасность сценариев оболочки}

Сценарии оболочки часто работают с файлами, внешним вводом и командами ОС, поэтому ошибки могут иметь серьёзные последствия.

\subsection{Типичные проблемы}

\begin{enumerate}
    \item отсутствие кавычек вокруг переменных;
    \item использование непроверенного пользовательского ввода;
    \item небезопасное создание временных файлов;
    \item выполнение строк как команд;
    \item предположение о фиксированном формате вывода внешних программ;
    \item зависимость от текущего каталога;
    \item неявная зависимость от \texttt{PATH}.
\end{enumerate}

\subsection{Пример ошибки из-за отсутствия кавычек}

Ненадёжный вариант:

\begin{lstlisting}
rm $filename
\end{lstlisting}

Если \texttt{filename} пустая или содержит пробелы, результат может быть неожиданным.

Лучше:

\begin{lstlisting}
rm -- "$filename"
\end{lstlisting}

Здесь:
\begin{enumerate}
    \item кавычки сохраняют значение как один аргумент;
    \item \texttt{--} обозначает конец опций и защищает от имён файлов, начинающихся с дефиса.
\end{enumerate}

\subsection{Строгий режим Bash}

В Bash часто используют:

\begin{lstlisting}
set -euo pipefail
\end{lstlisting}

Значение:
\begin{enumerate}
    \item \texttt{-e} --- завершать сценарий при ошибке команды;
    \item \texttt{-u} --- считать ошибкой обращение к неустановленной переменной;
    \item \texttt{pipefail} --- считать конвейер ошибочным, если ошиблась любая команда конвейера.
\end{enumerate}

Однако этот режим требует понимания исключений и особенностей поведения, особенно в условиях, циклах и конвейерах.

\section{Минимальные практические примеры}

\subsection{Резервное копирование каталога в shell}

\begin{lstlisting}
#!/bin/sh

src="$1"
dst="$2"

if [ "$#" -ne 2 ]; then
    echo "Использование: $0 SOURCE DESTINATION" >&2
    exit 1
fi

if [ ! -d "$src" ]; then
    echo "Источник не является каталогом" >&2
    exit 1
fi

mkdir -p "$dst"

archive="$dst/backup-$(date +%Y%m%d-%H%M%S).tar.gz"

tar -czf "$archive" "$src"

echo "Создан архив: $archive"
\end{lstlisting}

Шаги:
\begin{enumerate}
    \item проверить количество аргументов;
    \item проверить, что источник является каталогом;
    \item создать каталог назначения;
    \item сформировать имя архива с датой;
    \item вызвать \texttt{tar};
    \item вывести путь к архиву.
\end{enumerate}

\subsection{Поиск больших файлов в PowerShell}

\begin{lstlisting}
$path = "C:\Users"
$limit = 100MB

Get-ChildItem -Path $path -Recurse -File -ErrorAction SilentlyContinue |
    Where-Object { $_.Length -gt $limit } |
    Sort-Object Length -Descending |
    Select-Object FullName, Length
\end{lstlisting}

Шаги:
\begin{enumerate}
    \item получить файлы рекурсивно;
    \item отбросить ошибки доступа;
    \item оставить файлы больше заданного размера;
    \item отсортировать по размеру;
    \item вывести имя и размер.
\end{enumerate}

\subsection{Пакетный файл CMD для вывода файлов}

\begin{lstlisting}
@echo off

if "%~1"=="" (
    echo Использование: listfiles.bat DIRECTORY
    exit /b 1
)

if not exist "%~1" (
    echo Каталог не найден
    exit /b 1
)

for %%f in ("%~1\*") do (
    echo %%~nxf
)
\end{lstlisting}

Шаги:
\begin{enumerate}
    \item проверить наличие аргумента;
    \item проверить существование пути;
    \item пройти циклом по файлам;
    \item вывести имена файлов.
\end{enumerate}

\section{Типичные ошибки начинающих}

\subsection{Ошибки в shell}

\begin{enumerate}
    \item Пробелы вокруг \texttt{=}:
\begin{lstlisting}
x = 1   # ошибка
x=1     # правильно
\end{lstlisting}

    \item Отсутствие кавычек:
\begin{lstlisting}
cp $file $dest      # опасно
cp -- "$file" "$dest" # лучше
\end{lstlisting}

    \item Неверное сравнение строк и чисел:
\begin{lstlisting}
[ "$x" = "10" ]   # строковое сравнение
[ "$x" -eq 10 ]   # числовое сравнение
\end{lstlisting}

    \item Использование Bash-расширений в сценарии с \texttt{\#!/bin/sh}.
\end{enumerate}

\subsection{Ошибки в CMD}

\begin{enumerate}
    \item Использование \texttt{\%f} вместо \texttt{\%\%f} в пакетном файле.
    \item Непонимание поведения \texttt{if errorlevel N}.
    \item Проблемы с пробелами в путях без кавычек.
\end{enumerate}

\subsection{Ошибки в PowerShell}

\begin{enumerate}
    \item Ожидание текстового поведения конвейера вместо объектного.
    \item Непонимание различия между \texttt{\$?} и \texttt{\$LASTEXITCODE}.
    \item Использование строк там, где можно работать со свойствами объектов.
\end{enumerate}

\section{Краткие выводы}

Командная оболочка является важным компонентом пользовательского и административного взаимодействия с операционной системой. Она предоставляет интерактивную командную строку и язык сценариев для автоматизации.

В Unix-подобных системах доминируют оболочки семейства Bourne shell: \texttt{sh}, \texttt{bash}, \texttt{dash}, \texttt{ksh}, \texttt{zsh}. Их модель основана на запуске программ, текстовых потоках, кодах завершения, перенаправлениях и конвейерах.

В Windows традиционно используется \texttt{cmd.exe} с пакетными файлами \texttt{.bat} и \texttt{.cmd}. Для современной автоматизации Windows и кроссплатформенных задач применяется PowerShell, отличающийся объектным конвейером и тесной интеграцией с .NET.

Основные понятия, необходимые для уверенной работы с оболочками:
\begin{enumerate}
    \item команда;
    \item аргументы;
    \item переменные;
    \item переменные окружения;
    \item стандартные потоки;
    \item перенаправление;
    \item конвейер;
    \item код завершения;
    \item сценарий;
    \item встроенные и внешние команды;
    \item кавычки и экранирование;
    \item поиск команд через \texttt{PATH}.
\end{enumerate}

\end{document}
